\subsection{Características del Usuario}{\color{black}
\begingroup
\fuenteTabla\fontsize{10bp}{12bp}\selectfont
\setlength{\tabcolsep}{5bp}
\hyphenpenalty=10000 \exhyphenpenalty=10000 \sloppy
\renewcommand{\arraystretch}{1.2}
\begin{longtable}{|>{\raggedright\arraybackslash}p{0.28\linewidth}|>{\raggedright\arraybackslash}p{0.65\linewidth}|}
\caption{Características del usuario}\label{tab:caracteristicas-usuario}\\\hline
\textbf{Aspecto} & \textbf{Usuario único — bombero coordinador}\\\hline
\endfirsthead
\hline\textbf{Aspecto} & \textbf{Usuario único — bombero}\\\hline
\endhead
Características & Es el único tipo de usuario previsto e interactúa con la aplicación para llevar a cabo el caso de uso descrito en la sección 2.2.\\\hline
Nivel de privilegios & Todos los usuarios cuentan con un único nivel de privilegios y la aplicación no dispone de cuentas, autenticación ni modo administrador, como se indica en la sección 2.1.6.\\\hline
Rol & Interesado en visualizar la posible propagación del fuego.\\\hline
Estudios y experiencia técnica & Bomberos profesionales con experiencia en incendios y conocimientos básicos de física y química, familiarizados con propagación, convección, conducción, radiación y pirólisis. No se tienen conocimientos especializados en tecnología o simulación, ni saben utilizar o comprender los archivos resultantes de FDS.\\\hline
Frecuencia de uso & Cuando el usuario lo necesite, tanto en emergencias reales como en actividades de capacitación y estudio para bomberos en formación.\\\hline
\end{longtable}
\endgroup
}